\chapter{The data}
\label{ch:data}
\whereami{1}

\begin{tldr}
\begin{itemize}
\item 4,407 training exams, 24,371 series, one \textbf{DICOM} file per slice. About 1,300 hidden test exams.
\item Only \textbf{58} training exams have the 12 labels from doctors (``gold''). The other \textbf{4,349} have a
      free-text \textbf{radiology report} in one of 12 languages, and no labels.
\item The test exams have \textbf{images only}: no report.
\item So the real first task is: \textbf{turn 4,349 reports into labels}, then train an image model on them.
\end{itemize}
\end{tldr}

\section{Folders: study, series, slice}

\begin{figure}[H]
\centering
\begin{tikzpicture}[x=1mm,y=1mm, every node/.style={font=\footnotesize}]
  \node[anchor=west] at (0,0) {\faFolder\ \code{train\_series/}};
  \draw[muted] (2,-2) -- (2,-26);
  \draw[muted] (2,-6) -- (6,-6); \node[anchor=west] at (6,-6) {\faFolder\ \code{<StudyInstanceUID>/}};
  \node[anchor=west, text=us] at (62,-6) {\faArrowLeft\ one \textbf{study} = one exam = one row of labels};
  \draw[muted] (9,-8) -- (9,-24);
  \draw[muted] (9,-12) -- (13,-12); \node[anchor=west] at (13,-12) {\faFolder\ \code{<SeriesInstanceUID>/}};
  \node[anchor=west, text=us] at (62,-12) {\faArrowLeft\ one \textbf{series} = one 3D stack (e.g.\ sagittal PD fat-sat)};
  \draw[muted] (16,-14) -- (16,-21);
  \draw[muted] (16,-17) -- (20,-17); \node[anchor=west] at (20,-17) {\faFile\ \code{<SOPInstanceUID>.dcm}};
  \node[anchor=west, text=us] at (62,-17) {\faArrowLeft\ one \textbf{slice} (one 2D image + header)};
  \draw[muted] (16,-21) -- (20,-21); \node[anchor=west, text=muted] at (20,-21) {\ldots\ 20--40 more slices};
  \draw[muted] (9,-24) -- (13,-24); \node[anchor=west, text=muted] at (13,-24) {\ldots\ \ca 5 more series};
  \draw[muted] (2,-26) -- (6,-26); \node[anchor=west, text=muted] at (6,-26) {\ldots\ 4,406 more studies};
\end{tikzpicture}
\caption{The folder tree. The long numbers (UIDs) are random identifiers; they tell you nothing by themselves.
The CSV files link them to labels and to series descriptions.}
\end{figure}

\section{One DICOM file}

\term{DICOM} is the file format of medical imaging. A \code{.dcm} file has two parts:

\begin{figure}[H]
\centering
\begin{tikzpicture}[x=1mm,y=1mm]
  \node[draw=grid, fill=white, rounded corners=1.5mm, minimum width=84mm, minimum height=44mm] (f) at (0,0) {};
  \node[anchor=north west, font=\footnotesize\bfseries] at (-40,21) {header (``tags'')};
  \node[anchor=north west, font=\scriptsize\ttfamily, align=left, text=ink] at (-40,16)
    {Rows = 384, Columns = 384\\ PixelSpacing = 0.42 \textbackslash{} 0.42\\ SliceThickness = 3.5\\
     ImagePositionPatient = ...\\ ImageOrientationPatient = ...\\ Laterality = R \textcolor{muted}{(often missing)}\\
     \textcolor{muted}{... 86 allowed tags}};
  \fill[dark!88] (16,-17) rectangle (38,13);
  \fill[white!65!dark] (27,-2) circle (6.5);
  \fill[dark!88] (27,-2) circle (4);
  \node[font=\footnotesize\bfseries] at (27,17) {pixels};
  \node[font=\scriptsize, text=muted] at (27,-20) {16-bit numbers};
\end{tikzpicture}
\caption{A DICOM file = a dictionary of metadata + the image. The header matters as much as the pixels: it
says where the slice sits in the body, in millimetres.}
\end{figure}

\begin{center}\small
\begin{tabularx}{\linewidth}{l>{\raggedright\arraybackslash}X}
\toprule
Tag & Why we need it \\
\midrule
\code{ImagePositionPatient} & 3D position of the slice's corner (mm). Sort slices by this, never by file name. \\
\code{ImageOrientationPatient} & directions of the image rows and columns in the body. Tells the plane, and whether the image is flipped or rotated. \\
\code{PixelSpacing}, \code{SliceThickness} & size of a pixel and of a slice in mm. Needed to resize every series to the same physical scale. \\
\code{Laterality} & left or right knee. Missing in about half of the exams. \\
\bottomrule
\end{tabularx}
\end{center}

\begin{watchout}
The competition's files keep only 86 ``allowed'' header tags, and sites fill them inconsistently. Every tag can
be missing; read them defensively.
\end{watchout}

\section{The tables}

\begin{figure}[H]
\centering
\begin{tikzpicture}[x=1mm,y=1mm,
  tbl/.style={draw=grid, fill=white, rounded corners=1mm, align=left, font=\scriptsize, inner sep=2mm, anchor=north west}]
  \node[tbl, text width=44mm] at (0,0) {\textbf{train.csv}\\[1mm] one row per \textbf{study}\\
     \code{StudyInstanceUID}\\ 12 label columns (filled for the 58 gold studies only)\\ the \textbf{report} (free text)};
  \node[tbl, text width=44mm] at (52,0) {\textbf{train\_series.csv}\\[1mm] one row per \textbf{series}\\
     \code{SeriesInstanceUID}\\ \code{Anatomical\_Plane}: Sagittal / Coronal / Axial\\ \code{Fluid\_Sensitive}: 0 / 1\\
     \code{Fat\_Suppression}: 0 / 1};
  \node[tbl, text width=44mm] at (104,0) {\textbf{test\_series.csv}\\[1mm] the same for the hidden test\\
     series (only visible to the scoring run)\\[1mm] \textbf{sample\_submission.csv}\\ one row per test study,\\ 12 probability columns};
\end{tikzpicture}
\caption{The CSV files. \code{train\_series.csv} is precious: it already says which series is which, so we did
not have to guess from the images.}
\end{figure}

\section{The numbers}

\begin{figure}[H]
\centering
\begin{tikzpicture}[x=1mm,y=1mm,
  big/.style={font=\fontsize{20}{22}\selectfont\bfseries, text=accentdark},
  small/.style={font=\footnotesize, text=ink, align=center, text width=34mm}]
  \foreach \x/\n/\t in {0/{4,407}/{training exams}, 40/{24,371}/{series (\ca 5.5 per exam)}, 80/{\ca 1,300}/{hidden test exams},
                        120/12/{findings to predict}} {
    \node[big] at (\x,8) {\n};
    \node[small] at (\x,0) {\t};
  }
  \foreach \x/\n/\t/\c in {0/58/{exams with \textbf{doctor labels}}/them, 40/{4,349}/{exams with a \textbf{report} only}/us,
                           80/12/{report languages}/ink, 120/977/{characters: median report length}/ink} {
    \node[big, text=\c] at (\x,-16) {\n};
    \node[small] at (\x,-24) {\t};
  }
\end{tikzpicture}
\end{figure}

\section{Who labelled what}

\begin{figure}[H]
\centering
\begin{tikzpicture}[x=1mm,y=1mm]
  \foreach \i in {0,...,99} {
    \pgfmathsetmacro{\c}{mod(\i,25)}
    \pgfmathsetmacro{\r}{int(\i/25)}
    \fill[us!75] ({\c*4.6},{-\r*4.6}) rectangle ++(4,4);
  }
  % 58 of 4,407 = 1.3 squares of 44 studies each
  \fill[them] (0,0) rectangle ++(4,4);
  \fill[them] (4.6,0) rectangle ++(1.3,4);
  \draw[them, line width=0.8pt] (-1,5.2) -- (-1,6.5) -- (8,6.5) -- (8,5.2);
  \node[anchor=south west, font=\footnotesize, text=them!85!black] at (-1,6.6) {\textbf{58 gold exams} (1.3\,\%)};
  \node[anchor=west, font=\footnotesize, align=left] at (118,-6) {each square = 44 exams\\[1mm]
     \textcolor{them}{\rule{3mm}{3mm}} labelled by doctors\\ \hspace*{4mm}from the \textbf{images}\\[1mm]
     \textcolor{us!75}{\rule{3mm}{3mm}} only a \textbf{report}};
\end{tikzpicture}
\caption{The training set. Almost everything has to be learned from reports.}
\label{fig:waffle}
\end{figure}

\paragraph{The 58 gold exams.} Two musculoskeletal radiologists read the \emph{images} independently, a third
settled disagreements, and they applied the strict thresholds of Figure~\ref{fig:thresholds}. These are the only
labels we can fully trust.

\paragraph{The 4,349 report-only exams.} The report is what the hospital's radiologist wrote at the time,
in their own language and style. It mentions what that radiologist thought worth mentioning, with their own
severity words.

\begin{figure}[H]
\centering
\begin{tikzpicture}
\begin{axis}[xbar, width=11cm, height=6.2cm, bar width=2.6mm, xmin=0, xmax=1950,
    symbolic y coords={Thai,Serbian,Bulgarian,French,Dutch,Russian,German,Greek,Croatian,Turkish,Spanish,English},
    ytick=data, y tick label style={font=\footnotesize}, xtick=\empty, axis x line=none, axis y line*=left,
    nodes near coords, nodes near coords style={font=\scriptsize, anchor=west}, enlarge y limits=0.06]
  \addplot[barus] table[x=n, y=lang] {figdata/languages.dat};
\end{axis}
\end{tikzpicture}
\caption{Report languages of the 4,349 report-only exams. Any labelling method has to read all of them.}
\label{fig:languages}
\end{figure}

\begin{example}
A made-up report for our imaginary \textbf{Study K} (real reports are longer, about 1,000 characters):
\begin{quote}\small\itshape
Complete rupture of the anterior cruciate ligament. Bone contusion of the lateral femoral condyle and the
posterolateral tibial plateau. Moderate joint effusion. Menisci intact. Collateral ligaments intact.
Cartilage preserved.
\end{quote}
A careful reader would turn this into: ACL \textbf{yes}, Contusion \textbf{yes}, Effusion \textbf{yes}
(moderate), both menisci \textbf{no}, MCL \textbf{no}, all three OA \textbf{no}, and \ldots\ Synovitis, Baker's,
Fracture: \textbf{not mentioned}. What do you do with ``not mentioned''? Chapter~\ref{ch:labels} answers that.
\end{example}

\section{How common is each finding?}

\begin{figure}[H]
\centering
\begin{tikzpicture}
\begin{axis}[ybar, width=\linewidth, height=6cm, bar width=3.2mm, ymin=0, ymax=70,
    symbolic x coords={ACL,MCL,Med.~meniscus,Lat.~meniscus,Med.~OA,Lat.~OA,PF OA,Effusion,Synovitis,Baker's,Contusion,Fracture},
    xtick=data, x tick label style={font=\scriptsize, rotate=35, anchor=north east}, ylabel={\% of exams positive},
    ylabel style={font=\footnotesize}, ymajorgrids, legend style={at={(0.02,0.98)}, anchor=north west},
    enlarge x limits=0.05]
  \addplot[barthem] table[x=finding, y=gold] {figdata/prevalence.dat};
  \addplot[barus] table[x=finding, y=pool] {figdata/prevalence.dat};
  \legend{{58 gold exams (doctors)}, {4,349 report exams (our labels $>$ 0.5)}}
\end{axis}
\end{tikzpicture}
\caption{Each finding is about three times more common in the gold set than in the report set. The organisers
chose the gold exams to contain enough positives. Synovitis stands out: 47\,\% in gold, 10\,\% in our labels,
because reports rarely mention it.}
\label{fig:prevalence}
\end{figure}

\begin{watchout}
Because the gold set is enriched, a model's \emph{average} probabilities on gold look ``too low'' even when the
model is fine. AUC does not care (it only looks at order), but anything that relies on calibrated probabilities
would.
\end{watchout}

\section{The test set has no reports}

\begin{figure}[H]
\centering
\begin{tikzpicture}[x=1mm,y=1mm,
  box/.style={draw=grid, rounded corners=1.5mm, minimum height=13mm, align=center, font=\footnotesize, fill=white}]
  \node[font=\small\bfseries] at (-6,9) {training};
  \node[box, minimum width=34mm] at (30,9) {\faImage\ images};
  \node[box, minimum width=34mm, draw=us, line width=0.8pt] at (70,9) {\faFile\ report};
  \node[box, minimum width=34mm] at (110,9) {\faTag\ labels (58 only)};
  \node[font=\small\bfseries] at (-6,-9) {test};
  \node[box, minimum width=34mm] at (30,-9) {\faImage\ images};
  \node[box, minimum width=34mm, dashed, text=muted] at (70,-9) {no report};
  \node[box, minimum width=34mm, dashed, text=muted] at (110,-9) {no labels};
\end{tikzpicture}
\caption{At test time only the images exist. The report can teach the model during training, but the model has
to work from images alone.}
\end{figure}

\begin{keyidea}
Every solution in this competition has the same two steps:
\textbf{(1)} turn the reports (and maybe the images) of the 4,349 exams into labels;
\textbf{(2)} train an image model on those labels.
The image model can never be much better than its labels. Chapter~\ref{ch:labels} is step 1, Chapters
\ref{ch:cache}--\ref{ch:model} are step 2.
\end{keyidea}

\begin{wordcheck}
\term{Weak supervision}: training on labels that are cheaper but noisier than the real thing (here: labels read
from reports instead of doctors' image labels). \term{Gold labels}: the trusted reference labels (here: the 58).
\end{wordcheck}

\begin{checkpoint}
\begin{itemize}
\done I can describe the path from a folder of \code{.dcm} files to one study with its series.
\done I know which exams have labels, which have reports, and what the test set has.
\quiz Q4.1: Why sort slices by \code{ImagePositionPatient} rather than by file name?
\quiz Q4.2: Our labels say 11\,\% of exams have an ACL tear, the gold set says 41\,\%. Does that mean our labels
      are wrong?
\quiz Q4.3: Why can't the model simply read the report at test time?
\end{itemize}
\end{checkpoint}
